\section{Faithful flatness of actual Taylor extension and original analytic relations}

Write $R$ for the original holomorphic germ ring at the origin of
$\mathbb C^n$, $S=\mathbb C[[x_1,\ldots,x_n]]$, and $T:R\to S$ for
the genuine original Taylor homomorphism. All analytic coefficients below
belong to the original germ ring. Arbitrary formal coefficients are not
asserted to be convergent.

\begin{theorem}[The genuine original Taylor homomorphism is faithfully flat]
\label{hir:actual-original-taylor-faithfully-flat}
\uses{hir:actual-original-maximal-adic-formal-ring-equivalence,hir:actual-formal-extension-whole-original-ideal-contraction}
\lean{Hironaka.HironakaActualTaylorFaithfullyFlat.originalTaylorAlgHom_flat,Hironaka.HironakaActualTaylorFaithfullyFlat.originalTaylorAlgHom_faithfullyFlat}
\leanok
The actual homomorphism $T:R\to S$ is faithfully flat, including $n=0$.
\end{theorem}
\begin{proof}
The original holomorphic local ring is Noetherian. Its completion at the
actual maximal ideal is therefore flat over $R$. The previously constructed
actual completion isomorphism identifies this completed ring with $S$ and
identifies the composite original completion map with $T$ on every original
germ. An isomorphism is flat and the composite of flat homomorphisms is
flat, proving flatness for this specific genuine Taylor homomorphism.

For every original ideal $I$, the previously proved whole ideal contraction
identity gives $T^{-1}(IS)=I$. Thus, if $IS=S$, its contraction is $R$ and
$I=R$. Every proper ideal of $R$ consequently extends to a proper ideal of
$S$. Together with flatness, this is the proper-ideal criterion for faithful
flatness: the module $S$ is flat and $IS\ne S$ for every maximal ideal $I$.
No local-homomorphism assumption or proposed flatness certificate is input.
\end{proof}

\begin{theorem}[Every actual formal relation is generated by finite original analytic relations]
\label{hir:actual-formal-relations-original-analytic-syzygies}
\uses{hir:actual-original-taylor-faithfully-flat}
\lean{Hironaka.HironakaActualTaylorFaithfullyFlat.actual_formal_relation_iff_original_syzygies}
\leanok
Let $g_1,\ldots,g_r\in R$ and $A_1,\ldots,A_r\in S$ be any actual
finite original germ family and formal coefficient family. Then
\[
 \sum_i T(g_i)A_i=0
\]
if and only if there exist $k\geq0$, original germs $c_{ij}\in R$ and
formal series $B_j\in S$ such that
\[
 A_i=\sum_j T(c_{ij})B_j\quad\hbox{for every }i,
 \qquad \sum_i g_i c_{ij}=0\quad\hbox{in the original germ ring for every }j.
\]
The families may be empty. No analytic relation generators are supplied
as hypotheses.
\end{theorem}
\begin{proof}
Give $S$ the $R$-module structure induced by the actual homomorphism $T$.
The displayed vanishing relation is exactly $\sum_i g_i\mathbin{\cdot}A_i=0$
for that scalar action. Flatness implies the equational criterion for this
finite relation. Concretely, tensoring the map $R^r\to R$ defined by the
row $(g_i)$ preserves its kernel. The vector $(A_i)$ is in the resulting
kernel, so it is the image of an element of $S\otimes_R\ker(g_i)$.
Such a tensor is a finite sum $\sum_j B_j\otimes(c_{ij})_i$; its coordinates
give the first identities, and membership of each vector in the original
kernel gives the second identities. This is precisely the equational
criterion for flatness (Stacks Project, Tag 00HK).

Conversely, substitute the first identities into the relation and interchange
the two finite sums. Multiplicativity and additivity of $T$ give
$\sum_j T(\sum_i g_ic_{ij})B_j$, which is zero by the second identities.
This yields actual original analytic syzygies, while allowing the final
coefficients $B_j$ to remain formal.
\end{proof}
